\section{Introduction}
\makeatletter
\@afterindenttrue
\makeatother

Economic effects are rarely well summarized by a single average. The impact of a treatment may differ substantially across the outcome distribution, while the treatment itself may be endogenous and the empirical specification may require a large set of controls. Instrumental variable quantile regression (IVQR) provides a framework for studying such heterogeneous structural effects under endogeneity, and double/debiased machine learning (DML) extends this framework to settings with high-dimensional nuisance parameters. However, even with high-dimensional controls and orthogonalization, weak instruments can still remain a serious problem. When identification deteriorates, point estimates can become unstable and confidence regions can lose informativeness even when inference remains formally valid. In this thesis we study that problem directly by examining how progressively weaker instruments affect the finite-sample performance and informativeness of weak-identification-robust DML-IVQR inference, and whether these effects differ across the outcome distribution.

Quantile treatment effects describe how a structural effect varies across quantiles of the potential-outcome distribution. When treatment is endogenous, conventional quantile regression does not generally recover these effects because treatment choices may be related to unobserved determinants of the outcome. IVQR solves this problem by combining excluded instruments with restrictions on the structural ranks of potential outcomes \parencite{chernozhukov_hansen_2005}.

Two additional complications arise in modern applications. First, empirical specifications may contain many controls, transformations, and interactions. Their coefficients are nuisance parameters rather than the primary object of interest, but estimating a large nuisance model without structure can be unstable and may contaminate inference on the treatment effect. Regularization provides one way to estimate such high-dimensional components, while Neyman-orthogonal scores reduce the first-order sensitivity of the target moment to sufficiently small nuisance-estimation errors \parencite{chernozhukov_et_al_2018}. \textcite{chen_huang_tien_2021} combine these ideas with IVQR and develop a DML-IVQR procedure that handles high-dimensional controls while conducting candidate-specific inference on the structural quantile effect.

Second, regularization does not solve weak identification. Even a valid excluded instrument may contain little information about the endogenous treatment. When instrument-induced treatment variation is small, different candidate treatment-effect values can generate similar population moments, making them difficult to distinguish statistically. This problem is particularly important in IVQR because the relevant identifying restrictions are quantile specific and nonlinear. Standard first-stage measures indicate how strongly the instrument is related to treatment, but they do not fully capture how much information the instrument provides for identifying a structural quantile effect. The broader weak-instrument literature shows why conventional estimate-centered approximations can become unreliable when identifying information is weak, and why candidate-specific procedures can remain useful even when point identification is imprecise \parencite{staiger_stock_1997,stock_wright_2000,andrews_stock_sun_2019}. In IVQR, this distinction is reflected in weak-identification-robust test inversion: a procedure may remain asymptotically valid while producing broad sets of accepted candidate values and little power against economically relevant false alternatives.

High-dimensional nuisance estimation and weak instrument relevance are conceptually different problems. Neyman orthogonality reduces sensitivity to nuisance-estimation errors; it does not create identifying variation. Weak-identification-robust inference avoids relying on strong local identification, but cannot guarantee precise or informative confidence regions. Performance may therefore break down as instruments weaken even when nuisance estimation is handled well.

This distinction motivates the empirical question studied in this thesis. Existing Monte Carlo evidence for DML-IVQR evaluates regularized nuisance estimation in a high-dimensional design with a fixed level of instrument relevance. It therefore does not systematically show how the same estimators and inference procedures behave when excluded-instrument relevance is progressively weakened while the rest of the structural environment is held fixed. The central research question is: How does weakening instrument strength affect the finite-sample performance and informativeness of weak-identification-robust inference in DML-IVQR, and how do these effects vary across quantiles?

The thesis answers this question through a controlled extension of the Monte Carlo design in \textcite{chen_huang_tien_2021}. A single parameter, $\kappa \in \{1.00,0.50,0.25,0.10\}$, controls the loading of the excluded-instrument components in the latent treatment equation. Smaller values reduce excluded-instrument relevance. An independent component is added to the treatment index to preserve its marginal variance as $\kappa$ varies. The marginal treatment distribution, endogeneity, structural outcome equation, true quantile treatment effects, control dimension, and observed instrument construction are also preserved. This keeps the four relevance settings directly comparable while changing only how strongly the excluded instruments contribute to treatment. The parameter $\kappa$ is a controlled relevance parameter for this Monte Carlo design, not a universal scalar measure of IVQR identification strength.

Three estimators are compared on identical simulated samples. Oracle-GMM serves as an infeasible benchmark because it knows the true set of relevant controls. Full-GMM uses the full high-dimensional control vector without regularization. DML-IVQR uses regularized nuisance estimation together with the preserved Belloni--Chernozhukov penalty implementation. The comparison separates the common effect of declining instrument information from estimator-specific consequences of nuisance estimation, although differences between Full-GMM and DML-IVQR cannot be attributed to regularization alone because their residualization implementations also differ. Performance is evaluated across five quantiles, $\tau\in\{0.10,0.25,0.50,0.75,0.90\}$, and two sample sizes, $n\in\{500,1000\}$.

The evaluation separates point-estimation accuracy, coverage, power, and accepted-set informativeness. Bias, mean absolute error (MAE), and root mean squared error (RMSE) summarize accuracy. Coverage measures retention of the true structural quantile effect by the candidate-specific test, while power measures rejection of false treatment-effect values. Inferential informativeness is summarized by the grid-based accepted-set measure within the prespecified primary computational domain \(A_0=[-1,3]\); the choice of this domain and the corresponding domain-expansion diagnostic are explained in Appendix~\ref{app:domain-choice}. These distinctions matter: high coverage may accompany acceptance of many false candidates, while narrow accepted sets may reject the true value too frequently.

The Monte Carlo evidence shows a break down in informativeness as instrument relevance declines. Point-estimation error rises gradually for Oracle-GMM, Full-GMM, and DML-IVQR; accepted sets expand; and power against false candidate values falls. The deterioration of the oracle benchmark is particularly informative because it shows that the loss is not generated only by high-dimensional nuisance estimation. At the same time, the magnitude of the deterioration varies across quantiles and estimators. Coverage is also not uniformly close to its nominal level in finite samples, including in some strong-relevance cells. These results describe the implemented tests using a chi-square reference cutoff; the experiment does not establish the applicability of the asymptotic regularity conditions or identify the cause of undercoverage. DML-IVQR often performs better than the unregularized high-dimensional benchmark, but it does not dominate in every design cell and cannot compensate for the loss of identifying information when instruments become only weakly relevant. The main contribution of the thesis is therefore to document, within a controlled high-dimensional IVQR experiment, how nuisance robustness, instrument relevance, inferential informativeness, and finite-sample calibration interact as the information supplied by the excluded instruments is systematically reduced.
